\section*{Hoofdstuk \ref{ch:g5:biodiversity} --- Biodiversiteit}

\begin{solution}{exo:g5:biodiversity:1}
Eén bepaald type \omterm{def:g1:living-or-not:living}{levend wezen} --- individuen die op hetzelfde model
gebouwd zijn en zich onderling voortplanten. Voorbeelden: de huismus,
het madeliefje, de vos, de tuinslak.
\end{solution}

\begin{solution}{exo:g5:biodiversity:2}
De verscheidenheid aan leven die een plek herbergt --- vooral het aantal
\omterm{def:g5:biodiversity:species}{soorten} en hoe gezond die zijn. De bloemrijke wei herbergt tientallen
planten- en diersoorten op één vierkante meter: een hoge \omterm{def:g5:biodiversity:biodiversity}{biodiversiteit}.
Het bestrate plein herbergt er bijna geen.
\end{solution}

\begin{solution}{exo:g5:biodiversity:3}
Ongeveer twee miljoen. Omdat elke serieuze expeditie nog steeds
onbenoemde \omterm{def:g5:biodiversity:species}{soorten} vindt --- vooral kevers --- dus moeten er miljoenen
meer bestaan dan er in de catalogus staan.
\end{solution}

\begin{solution}{exo:g5:biodiversity:4}
Veel \omterm{def:g5:biodiversity:species}{soorten} in één \omterm{def:g2:where-animals-live:habitat}{leefgebied} (de vierkante meter wei); veel
\omterm{def:g2:where-animals-live:habitat}{leefgebieden} naast elkaar (vijver, heg, wei); verscheidenheid binnen één
\omterm{def:g5:biodiversity:species}{soort} (geen twee huisjes van tuinslakken helemaal gelijk).
\end{solution}

\begin{solution}{exo:g5:biodiversity:5}
In een rijk web kan een jager wiens prooi wegvalt op andere prooi
overstappen, en heeft elke rol invallers; de klap wordt opgevangen. In
een \omterm{def:g1:my-body:parts}{arm} web heeft elke \omterm{def:g5:biodiversity:species}{soort} weinig keuzes, dus sleept één verlies
andere mee.
\end{solution}

\begin{solution}{exo:g5:biodiversity:6}
Bijen en zweefvliegen: de bloesem bestuiven. Lieveheersbeestjes: de
bladluizen opeten. \omterm{prop:g3:sorting-living-things:groups}{Vogels}: op rupsen patrouilleren. Wormen: de bodem
open houden.
\end{solution}

\begin{solution}{exo:g5:biodiversity:7}
\omterm{def:g2:where-animals-live:habitat}{Leefgebied} vernield (de gedempte vijver); omgeving vergiftigd
(zwerfafval en chemische stoffen); er te veel van weggenomen
(overbevissing). \omterm{prop:g5:biodiversity:loss}{Uitgestorven} betekent volledig verdwenen --- overal,
voorgoed.
\end{solution}

\begin{solution}{exo:g5:biodiversity:8}
Open antwoord. Een begroeid of beplant vierkantje levert meestal tien of
meer \omterm{def:g5:biodiversity:species}{soorten} op; het bestrate vierkantje bijna nul --- het verschil is
de meting.
\end{solution}

\begin{solution}{exo:g5:biodiversity:9}
Elke \omterm{def:g5:biodiversity:species}{soort} is een draad: haar verlies laat elke buur --- alles wat haar
at, alles wat zij bestoof of opruimde --- met minder keuzes achter, en
verzwakt het vermogen van het web om de volgende klap op te vangen. En
omdat niemand alle draden kent die een \omterm{def:g5:biodiversity:species}{soort} vasthoudt, kan niemand een
verlies vooraf als onschadelijk verklaren; en omdat uitsterven voorgoed
is, zou die vergissing nooit ongedaan te maken zijn.
\end{solution}

\begin{solution}{exo:g5:biodiversity:10}
Ooievaars één voor één beschermen zou betekenen dat je \omterm{prop:g3:sorting-living-things:groups}{vogels} bewaakt;
natte weiden herstellen bouwde weer op waarvan de ooievaars
\emph{afhangen} --- en een \omterm{def:g2:where-animals-live:habitat}{leefgebied} voedt al zijn bewoners tegelijk,
en daarom keerden kikkers, libellen en weidebloemen met de ooievaars
terug. Geen \omterm{def:g2:where-animals-live:habitat}{leefgebied}, geen bewoners; \omterm{def:g2:where-animals-live:habitat}{leefgebied} hersteld, bewoners terug.
\end{solution}

\begin{solution}{exo:g5:biodiversity:11}
Wat weggehaald werd was geen versiering maar een personeelsbestand: de
\omterm{def:g5:biodiversity:species}{soorten} van de heg en de bloemenstrook huisvestten de bestuivers en de
plaageters, wier diensten in niemands boekhouding stonden zolang ze
gratis waren. Het werk van de natuur is onzichtbare boekhouding ---
voortdurend, onbetaald, onopgemerkt --- dus wordt de prijs ervan precies
ontdekt wanneer de werkers uitgezet worden en de taken als rekeningen
terugkomen.
\end{solution}
